\centerline{\bf \Huge Комбинаторика.} 
\centerline{\bf \Huge Правило сложения и умножения.}
\centerline{\bf \Huge Перестановки, размещения, сочетания.}
\title{Комбинаторика. \\ Правило сложения и умножения. \\ Перестановки, размещения, сочетания.}

1. Монету бросают 10 раз. Сколько разных последовательностей орлов и решек можно при этом получить?

2. В Калмыкии 4 известных населенных пункта: Элиста,Троицкое, Оргакин и Бага-Бурул. Из Элисты в Троицкое можно проехать по 4 дорогам, из Троицкого в Бага-Бурул по 5, из Элисты в Оргакин по 4, из Оргакина в Бага-Бурул по 6. Сколько всего способов добраться из Элисты в Бага-Бурул? 

3. Анаграммой называется произвольное слово, полученное из данного слова перестановкой букв. Сколько анаграмм можно составить из слов:

\par 

a) ОЧИР

б) МАГОМЕДРАСУЛ

в) ОЧИРМАГОМЕДРАСУЛ

г) ПРЕВОСОКОРАССМОТРИТЕЛЬСТВУЮЩИЙ

4. Из пяти борцов и двенадцати гимнастов нужно составить составить олимпийскую сборную из трех борцов и шести гимнастов. Сколькими способами можно составить сборную?

5. Сколькими способами можно выбрать из полной колоды (52 карты) 10 карт так, чтобы: \\
а) среди них был ровно один туз? \\ 
б) среди них был хотя бы один туз?
